% \setcounter{deso}{1}
\tatsode
\begin{name}
    {KỲ THI TUYỂN SINH VÀO LỚP 10}
    {ĐỀ CHÍNH THỨC}
    {$\star\star\star$}
    {Thời gian: 120 phút (không kể thời gian phát đề)}
\end{name}

\begin{bt}[1,5 điểm]~
    \begin{enumerate}[1.]
        \item Tính giá trị biểu thức $T=3\sqrt{2}+\sqrt{9}-\sqrt{18}$.
        \item Giải phương trình $2x-1=9$.
        \item Giải bất phương trình $2x-3\le 0$.
    \end{enumerate}
    \loigiai{
        \begin{enumerate}[1.]
            \item Ta có $\sqrt{9}=3$ và $\sqrt{18}=\sqrt{9\cdot 2}=3\sqrt{2}$. Do đó
                  \begin{center} $T=3\sqrt{2}+3-3\sqrt{2}=3$. \end{center}
            \item Giải phương trình $2x-1=9$.
                  \begin{eqnarray*}
                      2x&=&10\\
                      x&=&5
                  \end{eqnarray*}
            \item Giải bất phương trình $2x-3\le 0$.
                  \begin{eqnarray*}
                      2x&\le&3\\
                      x&\le&\dfrac32
                  \end{eqnarray*}
        \end{enumerate}
    }
\end{bt}

\begin{bt}[2,0 điểm]~
    \begin{enumerate}[1.]
        \item Cho biểu thức $M=\dfrac{\sqrt{x}}{\sqrt{x}+1}+\dfrac{2}{\sqrt{x}-1}-\dfrac{2}{x-1}$ với $x\ge 0$ và $ x\ne 1$.\\
              \begin{enumerate}
                  \item Rút gọn biểu thức $M$.
                  \item Tìm tất cả các giá trị của x để $M<\dfrac{1}{2}$.
              \end{enumerate}
        \item Gọi $x_1,x_2$ là hai nghiệm của phương trình $x^2-3x-6=0$. Không giải phương trình, hãy tìm số nguyên k nhỏ nhất thỏa mãn $2x_1^2-x_1x_2^2+7k\ge 0$.
    \end{enumerate}
    \loigiai{
        \begin{enumerate}[1.]
            \item Cho biểu thức $M=\dfrac{\sqrt{x}}{\sqrt{x}+1}+\dfrac{2}{\sqrt{x}-1}-\dfrac{2}{x-1}$ với $x\ge 0$ và $ x\ne 1$.\\
                  \begin{enumerate}
                      \item Với $x\ge 0$ và $x\ne 1$, ta có $x-1=(\sqrt{x}-1)(\sqrt{x}+1)$ và
                            \begin{eqnarray*}
                                M&=&\dfrac{\sqrt{x}}{\sqrt{x}+1}+\dfrac{2}{\sqrt{x}-1}-\dfrac{2}{(\sqrt{x}-1)(\sqrt{x}+1)}\\
                                &=&\dfrac{\sqrt{x}(\sqrt{x}-1)+2(\sqrt{x}+1)-2}{(\sqrt{x}-1)(\sqrt{x}+1)}\\
                                &=&\dfrac{x-\sqrt{x}+2\sqrt{x}+2-2}{x-1}\\
                                &=&\dfrac{x+\sqrt{x}}{x-1}\\
                                &=&\dfrac{\sqrt{x}(\sqrt{x}+1)}{(\sqrt{x}-1)(\sqrt{x}+1)}\\
                                &=&\dfrac{\sqrt{x}}{\sqrt{x}-1}.
                            \end{eqnarray*}
                      \item Ta có
                            \begin{eqnarray*}
                                M &<& \dfrac12\\
                                \dfrac{\sqrt{x}}{\sqrt{x}-1}-\dfrac12&<&0\\
                                \dfrac{\sqrt{x}+1}{2(\sqrt{x}-1)}&<&0.
                            \end{eqnarray*}
                            Vì tử thức $\sqrt{x}+1>0$ với mọi $x\ge 0$, nên dấu biểu thức phụ thuộc vào mẫu thức. Do đó
                            \begin{eqnarray*}
                                \sqrt{x}-1&<&0\\
                                \sqrt{x}&<&1 \\
                                0\le x &<&1
                            \end{eqnarray*}
                            Vậy $0\le x<1$ là giá trị cần tìm.
                  \end{enumerate}
            \item Với hai nghiệm $x_1,x_2$ của phương trình $x^2-3x-6=0$, theo hệ thức Viete ta có
                  \begin{eqnarray*}
                      x_1+x_2&=&3,\\
                      x_1x_2&=&-6.
                  \end{eqnarray*}
                  Do $x_1$ là nghiệm nên $x_1^2=3x_1+6$ và do $x_2$ là nghiệm nên $x_2^2=3x_2+6$. Khi đó
                  \begin{eqnarray*}
                      2x_1^2-x_1x_2^2+7k &=&2(3x_1+6)-x_1(3x_2+6)+7k\\
                      &=&6x_1+12-3x_1x_2-6x_1+7k\\
                      &=&6x_1+12-3(-6)-6x_1+7k\\
                      &=&30+7k.
                  \end{eqnarray*}
                  Do đó $2x_1^2-x_1x_2^2+7k\ge 0$ tương đương với
                  \begin{eqnarray*}
                      30+7k&\ge&0.
                  \end{eqnarray*}
                  Suy ra $k\ge -\dfrac{30}{7}$, nên số nguyên nhỏ nhất thỏa mãn là $k=-4$.
        \end{enumerate}
    }
\end{bt}

\begin{bt}[1,5 điểm]
    Theo kế hoạch, một đội xe thực hiện nhiệm vụ chở $120$ tấn lương thực đến vùng bị ảnh hưởng bởi thiên tai trong một số ngày quy định, mỗi ngày đội xe chở số tấn lương thực như nhau. Vì tính chất khẩn cấp nên đội xe đã được tăng cường thêm phương tiện vận chuyển. Do đó, mỗi ngày đội xe chở thêm được $10$ tấn lương thực so với kế hoạch. Nhờ vậy, đội xe đã hoàn thành nhiệm vụ sớm hơn $2$ ngày so với kế hoạch. Hỏi theo kế hoạch, mỗi ngày đội xe phải chở bao nhiêu tấn lương thực?
    \loigiai{
        Gọi $x$ (với $x>0$) là số tấn lương thực đội xe chở mỗi ngày theo kế hoạch. Khi đó số ngày chở lương thực theo kế hoạch là $\dfrac{120}{x}$ ngày.\\
        Vì tính chất khẩn cấp nên thực tế mỗi ngày đội xe chở $x+10$ tấn lương thực nên số ngày chở lương thực thực tế là $\dfrac{120}{x+10}$ ngày.\\
        Vì đội xe đã hoàn thành nhiệm vụ sớm hơn $2$ ngày so với kế hoạch nên
        \begin{eqnarray*}
            \dfrac{120}{x+10} &=& \dfrac{120}{x}-2.\\
            120x &=& 120(x+10)-2x(x+10)\\
            120x &=& 120x+1200-2x^2-20x\\
            2x^2+20x-1200 &=&0\\
            x^2+10x-600 &=&0
        \end{eqnarray*}
        Giải phương trình $x^2+10x-600 =0$ ta được hai nghiệm phân biệt $x_1=20$ và $x_2=-30$.\\ Vì $x>0$, nên $x=20$.\\
        Vậy theo kế hoạch, mỗi ngày đội xe phải chở $20$ tấn lương thực.
    }
\end{bt}

\begin{bt}[2,0 điểm]~
    \begin{enumerate}[1.]
        \item Tính thể tích và diện tích xung quanh của hình nón có chiều cao $h=12\,\text{cm}$ và bán kính đáy $r=5\,\text{cm}$.
        \item Xét phép thử: “Chọn ngẫu nhiên một số tự nhiên lớn hơn $10$ và nhỏ hơn $27$”.
              \begin{enumerate}
                  \item Viết không gian mẫu của phép thử đó.
                  \item Tính xác suất của biến cố $A$: “Số tự nhiên được chọn là số nguyên tố”.
              \end{enumerate}
    \end{enumerate}
    \loigiai{
        \begin{enumerate}[1.]
            \item Thể tích của hình nón là $V=\dfrac13 \pi r^2 h = \dfrac13 \cdot \pi \cdot 5^2 \cdot 12 = 100\pi $ (cm$^3$).\\
                  Độ dài đường sinh của hình nón là $l=\sqrt{r^2+h^2}=\sqrt{5^2+12^2}=13$ (cm).\\
                  Diện tích xung quanh của hình nón là $S_{xq}=\pi r l = \pi \cdot 5 \cdot 13 = 65\pi $ (cm$^2$).
            \item Xét phép thử: “Chọn ngẫu nhiên một số tự nhiên lớn hơn $10$ và nhỏ hơn $27$”.
                  \begin{enumerate}
                      \item Không gian mẫu của phép thử là
                            \begin{center}
                                $\Omega = \{11; 12; 13; 14;15;16;17;18;19;20;21;22;23;24;25;26\}$.
                            \end{center}
                      \item Số phần tử không gian mẫu là $n(\Omega)=16$.\\
                            Biến cố $A$: “Số tự nhiên được chọn là số nguyên tố” có tập các kết quả thuận lợi là
                            \begin{center} $\{13;17;19;23\}$. \end{center}
                            Suy ra số kết quả thuận lợi của biến cố $A$ là $n(A)=4$.\\
                            Vậy xác suất của biến cố $A$ bằng $P(A)=\dfrac{n(A)}{n(\Omega)}=\dfrac{4}{16}=\dfrac14$.
                  \end{enumerate}
        \end{enumerate}
    }
\end{bt}

\begin{bt}[3,0 điểm]
    Từ điểm M nằm bên ngoài đường tròn $(O;R)$, kẻ hai tiếp tuyến $MA$ và $MB$ với đường tròn $(O;R)$ ($A,B$ là các tiếp điểm).
    \begin{enumerate}[1.]
        \item Chứng minh tứ giác $AOBM$ nội tiếp.
        \item Kẻ đường kính $AC$ của đường tròn $(O;R)$, kẻ $BH\perp AC$ tại $H$. Gọi $K$ là giao điểm của $MC$ và $BH$, $D$ là giao điểm của hai đường thẳng $BC$ và $AM$. Chứng minh tam giác $MBD$ cân tại $M$ và $K$ là trung điểm của $BH$.
        \item Gọi $r$ là bán kính đường tròn nội tiếp tam giác $MAB$. Tính độ dài đoạn thẳng $BH$ theo $r$ và $R$.
    \end{enumerate}
    \loigiai{
        \begin{center}
            \begin{tikzpicture}[line join = round, line cap = round, thick, font = \small, scale = .7,declare function ={R=4;}]
                \path
                (0,0) coordinate (O)
                (7,0) coordinate (M);
                \path[name path global = circ1] (O) circle (R);
                \path[name path global = circ2] let \p1 = ($($(O)!.5!(M)$)-(M)$) in ($(O)!.5!(M)$) circle ({veclen(\x1,\y1)});
                \path[name intersections = {of = circ1 and circ2}]
                (intersection-1) coordinate (A)
                (intersection-2) coordinate (B)
                ;

                \path ($(O)!-1!(A)$) coordinate (C)
                ($(O)!(B)!(C)$) coordinate (H)
                (intersection of C--B and A--M) coordinate (D)
                (intersection of C--M and B--H) coordinate (K)
                ($(A)!.5!(B)$) coordinate (N);
                
                \path 
                (intersection of A--{$($(A)!1mm!(B)$)!.5!($(A)!1mm!(M)$)$} and B--{$($(B)!1mm!(A)$)!.5!($(B)!1mm!(M)$)$}) coordinate (I)
                ($(A)!(I)!(M)$) coordinate (P);

                \foreach \x/\y/\z in {C/B/A,D/A/O,M/B/O,B/H/A,I/P/M}{
                        \draw pic[draw, angle radius = 6pt]{right angle = \x--\y--\z};
                    }
                \draw (O)circle(4);
                \draw (A)--(C)--(D)--cycle (M)--(O)--(H)--(B)--cycle (M)--(C) (O)--(B)--(A)--(I)--(P);
                \foreach \d/\g in {O/135,A/45,B/-45,C/-135,M/45,H/135,D/-45,K/90,I/30,P/45}
                    {\fill (\d)circle(1.5pt) node[shift=(\g:10pt)]{$\d$};}
            \end{tikzpicture}
        \end{center}
        \begin{enumerate}[1.]
            \item Vì $MA$, $MB$ là hai tiếp tuyến của đường tròn $(O;R)$ ($A,B$ là các tiếp điểm) nên $MA \perp OA$, $MB \perp OB$.\\
                  Tam giác $MAO$ vuông tại $A$ có cạnh huyền $MO$ nên ba điểm $M$, $A$, $O$ cùng nằm trên đường tròn đường kính $MO$.\\
                  Tam giác $MBO$ vuông tại $B$ có cạnh huyền $MO$ nên ba điểm $M$, $B$, $O$ cùng nằm trên đường tròn đường kính $MO$.\\
                  Vậy bốn điểm $M$, $A$, $B$, $O$ cùng nằm trên đường tròn đường kính $MO$ nên tứ giác $AOBM$ nội tiếp đường tròn đường kính $MO$.
            \item Ta chứng minh tam giác $MBD$ cân tại $M$.\\
                  Trong đường tròn $(O;R)$ có $\widehat{ACB}$ là góc nội tiếp chắn cung $AB$, $\widehat{AOB}$ là góc ở tâm chắn cung $AB$ nên $\widehat{ACB}=\dfrac12 \widehat{AOB}$. \\
                  Vì $MA$, $MB$ là hai tiếp tuyến của đường tròn $(O;R)$ nên $OM$ là tia phân giác góc $\widehat{AOB}$.\\
                  Suy ra $\widehat{ACB}=\widehat{AOM}$. Do đó $OM \parallel BC$ (vì hai góc đồng vị).\\
                  Tam giác $ACD$ có $O$ là trung điểm cạnh $AC$, $OM \parallel CD$ nên có $OM$ là đường trung bình.\\
                  Khi đó $M$ là trung điểm đoạn $AD$ hay $MA=MD$.\\
                  Mà $MA=MB$ (tiếp tuyến đường tròn $(O;R)$) nên $MB=MD$.\\
                  Do vậy tam giác $MBD$ cân tại M.\\
                  Ta chứng minh $K$ là trung điểm của $BH$.\\
                  Vì $BH$ và $AD$ cùng vuông góc $AC$ nên $BH \parallel AD$.\\
                  Áp dụng định lí Thales trong tam giác $CAM$ ta được $\dfrac{HK}{AM}=\dfrac{CK}{CM}$.\\
                  Áp dụng định lí Thales trong tam giác $CDM$ ta được $\dfrac{BK}{DM}=\dfrac{CK}{CM}$.\\
                  Suy ra $\dfrac{HK}{AM}=\dfrac{BK}{DM}\left( =\dfrac{CK}{CM} \right) $.\\
                  Mà $AM=DM$ nên $HK=BK$. Do vậy $K$ là trung điểm $BH$.
            \item Gọi $J$ là giao điểm của tia $OM$ và đường tròn $(O;R)$.\\
                  Ta có $\widehat{BAJ}$ là góc nội tiếp chắn cung $BJ$, $\widehat{JOB}$ là góc ở tâm chắn cung $BJ$ nên $\widehat{BAJ}=\dfrac12 \widehat{JOB}$.\\
                  Mà $\widehat{JOB}=\widehat{JOA}$ (tính chất tiếp tuyến) và $\widehat{JOA}=\widehat{BAM}$ (cùng phụ $\widehat{OAB}$).\\
                  Do đó $\widehat{BAJ}=\dfrac12 \widehat{BAM}$ nên $AJ$ là tia phân giác góc $\widehat{MAB}$. Suy ra $J$ là giao điểm hai đường phân giác tam giác $MAB$.\\
                  Do đó $I$ trùng $J$ hay $I$ là giao điểm của $MO$ và đường tròn $(O;R)$.\\
                  Kẻ $IP \perp AM$. Khi đó $IP=r$ là bán kính đường tròn nội tiếp tam giác $MAB$.\\
                  Ta có $IP \parallel OA$ (vì cùng vuông góc $AM$). Áp dụng định lí Thales trong tam giác $MOA$ ta được
                  \begin{center}
                    $\dfrac{IP}{OA}=\dfrac{IM}{OM}=\dfrac{OM-OI}{OM} \Leftrightarrow \dfrac{r}{R}=\dfrac{OM-R}{OM}$
                  \end{center}
                  Suy ra 
                  \begin{center}
                    $OM \cdot r = OM \cdot R - R^2 \Leftrightarrow OM = \dfrac{R^2}{R-r}$. \quad ($*$)
                  \end{center}
                  Tam giác $BHC$ vuông tại $H$ có $BH=BC \cdot \sin ACB$.\\
                  Tam giác $ABC$ vuông tại $B$ (vì có $AC$ là đường kính đường tròn ngoại tiếp) có $BC=AC \cdot \cos ACB $.\\
                  Suy ra $BH=AC \cdot \cos ACB \cdot \sin ACB$.\\
                  Mà $\widehat{ACB}=\widehat{AOM}$ và trong tam giác $AOM$ vuông tại $A$ có 
                  \begin{itemize}
                  \item $\cos AOM = \dfrac{OA}{OM}= \dfrac{R-r}{R}$ (vì ($*$))\\
                  \item $\sin AOM = \dfrac{AM}{OM}= \sqrt{1-\cos ^2 AOM}=\sqrt{1-\dfrac{(R-r)^2}{R^2}}=\dfrac{\sqrt{2Rr-r^2}}{R}$.
                  \end{itemize}
                  Do vậy
                  \begin{center}
                    $BH = 2R \cdot \dfrac{R-r}{R} \cdot \dfrac{\sqrt{2Rr-r^2}}{R} = \dfrac{2}{R} \cdot (R-r) \cdot \sqrt{2Rr-r^2}$.
                  \end{center}
        \end{enumerate}
    }
\end{bt}